\documentclass[../../../include/open-logic-section]{subfiles}

\begin{document}

\olfileid{sth}{card-arithmetic}{ch}

\olsection{సాతత్య పరికల్పన}

అనంత కార్డినల్ సంఖ్యలు $2$ ఘాతాలతో, అంటే (\olref[opps]{lem:SizePowerset2Exp} ప్రకారం) ఘాత సమితులు తీసుకోవడంతో, ``నేరుగా'' ప్రవర్తిస్తాయని హామీ ఉంటే కార్డినల్ ఘాతాంకం \emph{సులభం} అవుతుందని ముందరి ఫలితం (సరిగ్గానే) సూచిస్తుంది. కానీ అనంత కార్డినల్ సంఖ్యలు ఘాత సమితులతో అలా ప్రవర్తిస్తాయని మనం అనుకోలేం.

ఈ విషయాన్ని విప్పి చెప్పడానికి ముందుగా చక్కని సంకేతాలను ప్రవేశపెడదాం.

\begin{defn}
$\cardsucc{\cardfont{a}}$ అనేది $\cardfont{a}$ కంటే కఠినంగా పెద్దదైన అత్యల్ప కార్డినల్ సంఖ్య అనుకుందాం. అప్పుడు రెండు అనంత శ్రేణులను ఇలా నిర్వచిస్తాం:
\begin{align*}
	\aleph_{0} &\defis \omega & 		
	\beth_{0} &\defis \omega\\
	\aleph_{\alpha \ordplus 1} &\defis \cardsucc{(\aleph_{\alpha})} &
	\beth_{\alpha+1} &\defis \cardexpo{2}{\beth_{\alpha}}\\
	\aleph_{\alpha} &\defis \bigcup_{\beta< \alpha} \aleph_{\beta} &
	\beth_{\alpha} &\defis \bigcup_{\beta < \alpha}\beth_{\beta} & \text{$\alpha$ సీమా క్రమసంఖ్య అయినప్పుడు}.
	\end{align*}
\end{defn}

$\cardsucc{\cardfont{a}}$ నిర్వచనం సమంజసం: ప్రతి కార్డినల్ సంఖ్య $\cardfont{a}$ కంటే పెద్ద కార్డినల్ సంఖ్య ఉందని \olref[cardinals][classing]{lem:NoLargestCardinal} చెబుతుంది; $\cardfont{a}$ కంటే పెద్దవాటిలో \emph{అత్యల్పమైనది} ఉందని అతిపరిమిత ఆగమనం నిర్ధారిస్తుంది. ఇక్కడ $\cardfont{a}$ స్థిర కార్డినల్ సంఖ్య; రెండు శ్రేణుల మిగిలిన నిర్వచనం అతిపరిమిత పునరావృత్తి ద్వారా వస్తుంది.
\sourcecorrection{OLTESTCARDCH-001}{మూలంలో ఈ పునరావృత్తి $\cardfont{a}$ నిర్వచనాన్ని పూర్తిచేస్తుందని ఉంది; అది స్థిర కార్డినల్ సంఖ్య కాదు, $\aleph$, $\beth$ శ్రేణుల నిర్వచనం. గద్యాన్ని సరిచేశాం.}

``$\aleph$'' సంకేతాన్ని కాంటర్ ప్రవేశపెట్టారు; ఇది హీబ్రూ వర్ణమాలలో మొదటి అక్షరమైన \emph{ఆలెఫ్}, ``అనంతం'' అనే హీబ్రూ పదంలోనూ మొదటి అక్షరం. ``$\beth$'' సంకేతాన్ని పియర్స్ ప్రవేశపెట్టారు; ఇది హీబ్రూ వర్ణమాలలో రెండవ అక్షరమైన \emph{బెత్}.\footnote{1900 డిసెంబరులో కాంటర్‌కు రాసిన లేఖలో పియర్స్ ఈ సంకేతాన్ని వాడారు. దురదృష్టవశాత్తూ, $\alpha \geq
\omega$కు $\beth_\alpha$ లేదనే తప్పు వాదననూ అక్కడ ఇచ్చారు.} ఈ సంకేతాలు ఇప్పుడు మనకు అనంత కార్డినల్ సంఖ్యలను ఇస్తాయి.

\begin{prop}
ప్రతి క్రమసంఖ్య $\alpha$కూ $\aleph_\alpha$, $\beth_\alpha$ కార్డినల్ సంఖ్యలు.
\end{prop}

\begin{proof}
రెండు ఫలితాలూ సరళమైన అతిపరిమిత ఆగమనంతో వస్తాయి. \olref[cardinals][classing]{omegaisacardinal} ప్రకారం $\aleph_0 =
\beth_0 = \omega$ కార్డినల్ సంఖ్య. $\aleph_\alpha$, $\beth_\alpha$ రెండూ కార్డినల్ సంఖ్యలని అనుకుంటే, $\aleph_{\alpha+1}$, $\beth_{\alpha+1}$ నిర్వచనాల ప్రకారమే కార్డినల్ సంఖ్యలు. \olref[cardinals][classing]{unioncardinalscardinal} ప్రకారం కార్డినల్ సంఖ్యల సమితి సమ్మేళనం కూడా కార్డినల్ సంఖ్య.
\end{proof}
\noindent
అంతేకాక ప్రతి అనంత కార్డినల్ సంఖ్య ఒక $\aleph$:

\begin{prop}
$\cardfont{a}$ అనంత కార్డినల్ సంఖ్య అయితే, ఏకైక $\gamma$కు $\cardfont{a} =
\aleph_\gamma$.
\end{prop}

\begin{proof}
అనంత కార్డినల్ సంఖ్యలపై అతిపరిమిత ఆగమనం చేద్దాం. ఆధార సందర్భం మొదటి ఆలెఫ్ నిర్వచనమే. ఆగమనంలో $\cardfont{b} < \cardfont{a}$ అయ్యే అనంత కార్డినల్ సంఖ్యకు $\cardfont{b} =
\aleph_{\gamma_\cardfont{b}}$ అని అనుకుందాం. ఏదో అనంత $\cardfont{b}$కు $\cardfont{a} =
\cardsucc{\cardfont{b}}$ అయితే $\cardfont{a} =
\cardsucc{(\aleph_{\gamma_\cardfont{b}})}=
\aleph_{\gamma_\cardfont{b}+1}$. $\cardfont{a}$ ఏ అనంత కార్డినల్ సంఖ్యకూ ఉత్తరవర్తి కాకపోతే, కార్డినల్ సంఖ్యలూ క్రమసంఖ్యలే కాబట్టి $\cardfont{a} = \bigcup_{\cardfont{b} < \cardfont{a}} \cardfont{b} =
\bigcup_{\omega\leq\cardfont{b} < \cardfont{a}}{\aleph_{\gamma_\cardfont{b}}}$; అందువల్ల $\gamma =
\bigcup_{\omega\leq\cardfont{b} < \cardfont{a}}\gamma_\cardfont{b}$ అనుకుంటే $\cardfont{a} = \aleph_\gamma$. శ్రేణి ప్రతి ఉత్తరవర్తి దశలో కఠినంగా పెరుగుతుంది; సీమా దశలో పూర్వ విలువల సమ్మేళనం తీసుకుంటుంది. కాబట్టి సూచిక ఏకైకం.
\sourcecorrection{OLTESTCARDCH-002}{మూల ఆగమన పరికల్పన పరిమిత కార్డినల్ సంఖ్యలనూ ఆలెఫ్ రూపంలోకి తీసుకుంది; సీమా సమ్మేళనంలో వాటికి నిర్వచించని సూచికలను వాడింది. ఆధారంగా $\omega$ను చూపి, ఆలెఫ్-సూచిక కలిగిన అనంత పూర్వ కార్డినల్ సంఖ్యలకే రెండవ సమ్మేళనం, సూచిక సమ్మేళనం పరిమితం చేశాం; ఏకైకత్వ కారణాన్ని జోడించాం.}
\end{proof}

ప్రతి అనంత కార్డినల్ సంఖ్య $\aleph$ రూపంలో ఉంటుంది. మరి ప్రతి అనంత కార్డినల్ సంఖ్య $\beth$ రూపంలోనూ ఉంటుందా? అలా \emph{ఉంటే}, అనంత కార్డినల్ సంఖ్యలు ఘాత సమితి క్రియతో ``నేరుగా'' ప్రవర్తిస్తాయి. నిజానికి అప్పుడు ఈ కింది వాదన మనకు ఉంటుంది:

\begin{defish}
\emph{సాధారణీకృత సాతత్య పరికల్పన} (GCH). అన్ని $\alpha$లకు $\aleph_\alpha  = \beth_\alpha$.
\end{defish}

ఇంకా, GCH నిజమైతే కార్డినల్ ఘాతాంకాన్ని గణనీయంగా సరళీకరించవచ్చు. ముఖ్యంగా $\cardfont{a}$ అనంతమై, చిన్న ఘాతం శూన్యం కాని కార్డినల్ సంఖ్యగా $\cardfont{b} < \cardfont{a}$ అయినప్పుడు $\cardexpo{\cardfont{a}}{\cardfont{b}}$ విలువ $\cardfont{a}\leq \cardexpo{\cardfont{a}}{\cardfont{b}} \leq
\cardsucc{\cardfont{a}}$ మధ్య ఉంటుందని చూపవచ్చు. ఆ రెండు అవకాశాల్లో ఏది వస్తుందో నిర్ణయించే ఖచ్చితమైన షరతులనూ చెప్పవచ్చు (అంటే $\cardfont{a} = \cardexpo{\cardfont{a}}{\cardfont{b}}$నా, లేక $\cardexpo{\cardfont{a}}{\cardfont{b}} =
\cardsucc{\cardfont{a}}$నా).\footnote{ఆ షరతును \emph{సహాంత్యత} నిర్ణయిస్తుంది.}
\sourcecorrection{OLTESTCARDCH-003}{మూల గరిష్ఠ-సరిహద్దు వాదనలో ఘాతం శూన్యం కావచ్చని అనుమతించారు; అప్పుడు అనంత ఆధారానికి ఘాత ఫలితం ఒకటి, ఆధారానికి కనీసం సమానం కాదు. అనంత ఆధారం, శూన్యం కాని చిన్న ఘాతం అని అర్హత చేర్చాం.}

కానీ GCH ఒక \emph{పరికల్పన}, \emph{సిద్ధాంతం} కాదు. నిజానికి $\ZFC$ సుసంగతమైతే $\ZFC + \text{GCH}$ కూడా సుసంగతమని \citet{Godel1938} నిరూపించారు. తరువాత $\lnot$GCHనూ $\ZFC$కు చేర్చవచ్చని తేలింది. ముఖ్యంగా GCHలోని అతి సరళమైన ప్రాధాన్యమైన \emph{ఉదాహరణ}ను చూడండి:

\begin{defish}
\emph{సాతత్య పరికల్పన} (CH). $\aleph_1 = \beth_1$.
\end{defish}

$\ZFC$ సుసంగతమైతే $\ZFC
+ \lnot\text{CH}$ కూడా సుసంగతమని \citet{Cohen1963} నిరూపించారు. కాబట్టి సాతత్య పరికల్పన $\ZFC$ నుంచి స్వతంత్రం.

``సాతత్యం'' అనేది వాస్తవ సంఖ్యారేఖ $\Real$కు మరొక పేరు కాబట్టి దీనిని సాతత్య పరికల్పన అంటారు. $\card{\Real} =
\beth_1$ అని \olref[opps]{continuumis2aleph0} చెబుతుంది. అందువల్ల సహజ సంఖ్యల కార్డినాలిటీ $\aleph_0 = \beth_0$కూ, సాతత్యం కార్డినాలిటీ $\beth_1$కూ మధ్య కార్డినల్ సంఖ్య ఏదీ లేదని సాతత్య పరికల్పన చెబుతుంది.

$\ZFC$ నుంచి (G)CH \emph{స్వతంత్రం} అని తెలిసిన తరువాత, వాటి \emph{సత్యం} గురించి ఏం చెప్పాలి? చెప్పాల్సింది \emph{చాలా} ఉంది. నిజానికి సమితి సిద్ధాంతంలో ఇటీవలి ఫలప్రదమైన పరిశోధనలో చాలా భాగం ఈ ప్రశ్నలనే పరిశీలిస్తోంది. అయితే రెండు చిన్న విషయాలను ప్రత్యేకంగా చెప్పాలి.

మొదటిది: ఈ ఔపచారిక స్వతంత్రత ఫలితాల నుంచి GCH లేదా CH సత్యవిలువ \emph{అనిర్ణీతం} అని \emph{వెంటనే} తేలదు. ప్రశ్నను ఒక వైపో మరొక వైపో నిర్ణయించే, మనకు సహజంగా అనిపించే మరిన్ని స్వయంసిద్ధాలను చేర్చాల్సి ఉండవచ్చు. గోడెల్ స్వయంగా ఇదే సరైన ప్రతిస్పందన అని సూచించారు.

రెండవది: $\ZFC$ నుంచి CH స్వతంత్రత నిశ్చయంగా \emph{ఆశ్చర్యకరం}; కానీ అక్షరార్థంలో \emph{నమ్మశక్యం కానిది} మాత్రం కాదు. కారణం సులభం: $\ZFC$ చెప్పినంతవరకు, ఒక కార్డినల్ సంఖ్య నుంచి దాని ఉత్తరవర్తికి వెళ్లడానికి ఘాత సమితి తీసుకోవడంకంటే సూక్ష్మమైన పద్ధతి అవసరం కావచ్చు.
 %The operation of taking powersets moves us from one stage of the
 %hierarchy to its successor stage (from $V_{\alpha}$ to
 %$V_{\alpha+1}$). 

ఈ రెండు విషయాల తరువాత మరింత తెలుసుకోవాలనుకుంటే, ఈ వాదంలో తమ తమ అభిప్రాయాలున్న తత్త్వవేత్తలు, గణితవేత్తల రచనలవైపు వెళ్లాలి.\footnote{మొదట \citet[\S15.6]{Potter2004} చదవడం మంచిదేమో.}

\end{document}
